\renewcommand{\Anno}{1997}

\renewcommand{\NumeroProblema}{5}

\renewcommand{\Testo}{
Costruire (con riga e compasso) un triangolo $ABC$ di cui sono dati la lunghezza $2p$ del perimetro, l'ampiezza $\alpha$ dell'angolo in $A$ e la lunghezza $h$ dell'altezza relativa al lato $BC$.
}

\renewcommand{\Risposta}{
% SSSUP-BOZZA-AI: Risposta; autore=ChatGPT; data=2026-11-08; stato=DA-VERIFICARE
% SSSUP-PROVENIENZA: AI: sintesi della soluzione preesistente
Per $p,h>0$ e $0<\alpha<\pi$, la costruzione è possibile se e solo se
\[
h\cos(\alpha/2)\leq p\bigl(1-\sin(\alpha/2)\bigr).
\]
Si disegna l'angolo di vertice $A$ e ampiezza $\alpha$, si prende $S$ su una sua semiretta con $AS=p$ e si costruisce la circonferenza tangente alle due semirette in $S$ e nel punto simmetrico. Una tangente comune interna a questa circonferenza e a quella di centro $A$ e raggio $h$ individua, sulle due semirette, i vertici $B,C$.
}

\renewcommand{\Aiutino}{
Disegniamo l'angolo con vertice in $A$.  La retta $BC$ deve essere
\ldots\ alla circonferenza di centro $A$ e raggio $h$.  Fissato, su
una delle semirette che delimitano l'angolo, un punto $T$ tale che $AT=p$, allora $T$ è il punto di tangenza tra una circonferenza ex-inscritta ad $ABC$ e la semiretta in questione \ldots.  Il problema si riconduce quindi alla costruzione delle bitangenti a due circonferenze date.  Per risolvere quest'ultimo problema, conviene intanto tracciare il punto in cui le bitangenti incontrano la congiungente i centri.
}

\renewcommand{\SoluzioneUfficiale}{\TitoloSoluzione Iniziamo
ricordando alcune proprietà dei triangoli che ci serviranno nella costruzione.  Sia $ABC$ un triangolo, sia $H$ il piede dell'altezza uscente da $A$, sia $O$ il centro della circonferenza ex-inscritta tangente al lato $BC$ in $K$ ed ai prolungamenti di $AB$ e $BC$, rispettivamente, in $S$ e $T$.  Allora $BC$ è tangente alla circonferenza con centro in $A$ e raggio $AH$ (in quanto $BC$ è perpendicolare al raggio $AH$) ed inoltre $AS=AT=p$.  Per dimostrare quest'ultima uguaglianza (che è una ben nota proprietà delle circonferenze ex-inscritte) basta osservare che $BK=BS$, $CK=CT$, per cui
\begin{eqnarray*}
	2p & = & AB+BC+AC  \\
	 & = & AB+BK+KC+AC  \\
	 & = & AB+BS+CT+AC  \\
	 & = & AS+AT,
\end{eqnarray*}
e concludere osservando che $AS=AT$.

\begin{center}
	\psset{unit=15mm}
	\pspicture(0,-1.6)(6,3)
	\rput(3,-1.2){Figura 1}
	
	\pstGeonode[PosAngle={-90,0,-90}](0,0){A}(2.5,1.5){B}(4,0){C}
	\pstProjection{B}{C}{A}[H]
	\pstHomO[HomCoef=2,PointSymbol=none]{A}{B}[B'']
	\pstRotation[AngleCoef=0.5,RotAngle=\pstAngleAOB{C}{A}{B},PointSymbol=none]{A}{C}[A']
	\pstRotation[AngleCoef=0.5,RotAngle=\pstAngleAOB{C}{B}{B''},PointSymbol=none]{B}{C}[B']
	\pstInterLL{A}{A'}{B}{B'}{O}
	\pstProjection{A}{B}{O}[S]
	\pstProjection{A}{C}{O}[T]
	\pstProjection[PosAngle=45]{B}{C}{O}[K]
	\pstRotation[AngleCoef=1.8,RotAngle=\pstAngleAOB{C}{A}{H},PointSymbol=none]{A}{C}[H']
	\pstRotation[AngleCoef=1.3,RotAngle=\pstAngleAOB{H}{A}{A},PointSymbol=none]{A}{H}[C']
	\pstArcOAB[linecolor=green]{A}{C'}{H'}
	
	\pstCircleOA[linecolor=green]{O}{S}
	\pstLineAB[nodesepB=-1.2]{A}{S}
	\pstLineAB[nodesepB=-1.2]{A}{T}
	\pstLineAB[nodesepB=-1.2]{C}{H}
	
	\pstLineAB[linewidth=1.5\pslinewidth]{A}{B}
	\pstLineAB[linewidth=1.5\pslinewidth]{A}{C}
	\pstLineAB[linewidth=1.5\pslinewidth]{C}{B}
	
	\pstLineAB[linecolor=blue,linestyle=dashed]{A}{H}
	\pstRightAngle[linecolor=blue]{A}{H}{B}
	
	\endpspicture
\end{center}


Sfruttando tali proprietà possiamo ottenere il triangolo richiesto mediante la seguente costruzione.
\begin{itemize}
	\item Supponiamo dato il punto $A$ e 2 semirette uscenti da $A$ che delimitano un angolo di ampiezza $\alpha$, e tracciamo la circonferenza con centro in $A$ e raggio $h$ (Figura 2).

	\item Segniamo su una delle 2 semirette un punto $S$ tale che $AS=p$ e tracciamo la circonferenza passante per $P$ e tangente alle due semirette (il suo centro è il punto d'incontro tra la bisettice dell'angolo in $A$ e la perpendicolare ad $AS$ passante per $S$, vedi Figura 3).

	\item  Tracciamo una delle 2 rette tangenti ad entrambe le circonferenze e che lasciano i centri in semipiani diversi. Le intersezioni tra questa retta e le 2 semirette iniziali sono i rimanenti vertici del triangolo desiderato (Figura 4).
\end{itemize}

\noindent
\psset{unit=10mm}
\pspicture(-1,-2)(6.8,4)
\rput(2.9,-1.6){Figura 2}
\psframe[linewidth=.6\pslinewidth](-1,-1.2)(6.8,4)

\pstGeonode[PosAngle={-90,0,-90},PointSymbol={*,none}](0,0){A}(2.5,1.5){B}(4,0){C}
\pstProjection[PointSymbol=none]{B}{C}{A}[H]
\pstHomO[HomCoef=2,PointSymbol=none]{A}{B}[B'']
\pstRotation[AngleCoef=0.5,RotAngle=\pstAngleAOB{C}{A}{B},PointSymbol=none]{A}{C}[A']
\pstRotation[AngleCoef=0.5,RotAngle=\pstAngleAOB{C}{B}{B''},PointSymbol=none]{B}{C}[B']
\pstInterLL[PointSymbol=none]{A}{A'}{B}{B'}{O}
\pstProjection[PointSymbol=none]{A}{B}{O}[S]
\pstProjection[PointSymbol=none]{A}{C}{O}[T]
\pstProjection[PosAngle=45,PointSymbol=none]{B}{C}{O}[K]
\pstRotation[AngleCoef=1.8,RotAngle=\pstAngleAOB{C}{A}{H},PointSymbol=none]{A}{C}[H']
\pstRotation[AngleCoef=1.3,RotAngle=\pstAngleAOB{H}{A}{A},PointSymbol=none]{A}{H}[C']
\pstArcOAB[linecolor=green]{A}{C'}{H'}
\pstMarkAngle[MarkAngleRadius=0]{C}{A}{B}{$\alpha$}
\pstRotation[RotAngle=15,PointSymbol=none]{A}{H}[H'']
\pstLineAB[linecolor=blue,linestyle=dashed]{A}{H''}
\Aput{$h$}

%\pstCircleOA[linecolor=green]{O}{S}
\pstLineAB[nodesepB=-1.2]{A}{S}
\pstLineAB[nodesepB=-1.2]{A}{T}
%\pstLineAB[nodesepB=-1.2]{C}{H}

%\pstLineAB[linewidth=1.5\pslinewidth]{A}{B}
%\pstLineAB[linewidth=1.5\pslinewidth]{A}{C}
%\pstLineAB[linewidth=1.5\pslinewidth]{C}{B}

%\pstLineAB[linecolor=blue,linestyle=dashed]{A}{H}
%\pstRightAngle[linecolor=blue]{A}{H}{B}

\endpspicture
\hfill\psset{unit=10mm}
\pspicture(-1,-2)(6.8,4)
\rput(2.9,-1.6){Figura 3}
\psframe[linewidth=.6\pslinewidth](-1,-1.2)(6.8,4)

\pstGeonode[PosAngle={-90,90,-90},PointSymbol={*,none}](0,0){A}(2.5,1.5){B}(4,0){C}
\pstProjection[PointSymbol=none]{B}{C}{A}[H]
\pstHomO[HomCoef=2,PointSymbol=none]{A}{B}[B'']
\pstRotation[AngleCoef=0.5,RotAngle=\pstAngleAOB{C}{A}{B},PointSymbol=none]{A}{C}[A']
\pstRotation[AngleCoef=0.5,RotAngle=\pstAngleAOB{C}{B}{B''},PointSymbol=none]{B}{C}[B']
\pstInterLL[PosAngle=-20]{A}{A'}{B}{B'}{O}
\pstProjection[PosAngle=80]{A}{B}{O}[S]
\pstProjection[PointSymbol=none]{A}{C}{O}[T]
\pstProjection[PosAngle=45,PointSymbol=none]{B}{C}{O}[K]
\pstRotation[AngleCoef=1.8,RotAngle=\pstAngleAOB{C}{A}{H},PointSymbol=none]{A}{C}[H']
\pstRotation[AngleCoef=1.3,RotAngle=\pstAngleAOB{H}{A}{A},PointSymbol=none]{A}{H}[C']
\pstArcOAB[linecolor=green]{A}{C'}{H'}
\pstMarkAngle[MarkAngleRadius=0.5,linecolor=blue]{C}{A}{O}{}
\pstMarkAngle[MarkAngleRadius=0.6,linecolor=blue]{C}{A}{O}{}
\pstMarkAngle[MarkAngleRadius=0.7,linecolor=blue]{O}{A}{B}{}
\pstMarkAngle[MarkAngleRadius=0.8,linecolor=blue]{O}{A}{B}{}

\pstCircleOA[linecolor=green]{O}{S}
\pstLineAB[nodesepB=-1.2]{A}{S}
\pcline[offset=9pt,linecolor=blue,linewidth=.6\pslinewidth]{|-|}(A)(S)
\lput*{:U}{$p$}
\pstLineAB[nodesepB=-1.2]{A}{T}
\pstLineAB[nodesepB=-0.5,linecolor=blue,linestyle=dashed]{A}{O}
\pstLineAB[nodesepB=-0.5,linecolor=blue,linestyle=dashed]{S}{O}
\pstRightAngle[linecolor=blue]{A}{S}{O}
%\pstLineAB[nodesepB=-1.2]{C}{H}

%\pstLineAB[linewidth=1.5\pslinewidth]{A}{B}
%\pstLineAB[linewidth=1.5\pslinewidth]{A}{C}
%\pstLineAB[linewidth=1.5\pslinewidth]{C}{B}

%\pstLineAB[linecolor=blue,linestyle=dashed]{A}{H}
%\pstRightAngle[linecolor=blue]{A}{H}{B}

\endpspicture

\begin{center}
	\psset{unit=10mm}
	\pspicture(-1,-2)(6.8,4)
	\rput(2.9,-1.6){Figura 4}
	\psframe[linewidth=.6\pslinewidth](-1,-1.2)(6.8,4)
	
	\pstGeonode[PosAngle={-90,90,-90}](0,0){A}(2.5,1.5){B}(4,0){C}
	\pstProjection[PointSymbol=none]{B}{C}{A}[H]
	\pstHomO[HomCoef=2,PointSymbol=none]{A}{B}[B'']
	\pstRotation[AngleCoef=0.5,RotAngle=\pstAngleAOB{C}{A}{B},PointSymbol=none]{A}{C}[A']
	\pstRotation[AngleCoef=0.5,RotAngle=\pstAngleAOB{C}{B}{B''},PointSymbol=none]{B}{C}[B']
	\pstInterLL{A}{A'}{B}{B'}{O}
	\pstProjection{A}{B}{O}[S]
	\pstProjection[PointSymbol=none]{A}{C}{O}[T]
	\pstProjection[PosAngle=45,PointSymbol=none]{B}{C}{O}[K]
	\pstRotation[AngleCoef=1.8,RotAngle=\pstAngleAOB{C}{A}{H},PointSymbol=none]{A}{C}[H']
	\pstRotation[AngleCoef=1.3,RotAngle=\pstAngleAOB{H}{A}{A},PointSymbol=none]{A}{H}[C']
	\pstArcOAB[linecolor=green]{A}{C'}{H'}
	
	\pstCircleOA[linecolor=green]{O}{S}
	\pstLineAB[nodesepB=-1.2]{A}{S}
	\pstLineAB[nodesepB=-1.2]{A}{T}
	\pstLineAB[nodesepA=-1.2,nodesepB=-1.2]{C}{H}
	
	\pstLineAB[linewidth=1.5\pslinewidth]{A}{B}
	\pstLineAB[linewidth=1.5\pslinewidth]{A}{C}
	\pstLineAB[linewidth=1.5\pslinewidth]{C}{B}
	
	%\pstLineAB[linecolor=blue,linestyle=dashed]{A}{H}
	%\pstRightAngle[linecolor=blue]{A}{H}{B}
	
	\endpspicture
\end{center}

La condizione per l'esistenza delle rette bitangenti che ci interessano è che le 2 circonferenze non siano secanti (se sono tangenti, la retta tangente comune risolve il problema), e cioè che la distanza tra i loro centri sia maggiore od uguale della somma dei raggi. Nel nostro caso, con un po' di trigonometria elementare si ricava che
$$AO=\frac{AS}{\cos(\alpha/2)}=\frac{p}{\cos(\alpha/2)},
\hspace{3em}
AH=h,
\hspace{3em}
OS=AS\cdot\tan(\alpha/2)=p\tan(\alpha/2),$$
e pertanto la condizione diventa che
$$\frac{p}{\cos(\alpha/2)}\geq h+p\tan(\alpha/2),$$
o equivalentemente
$$h\cos(\alpha/2)\leq p(1-\sin(\alpha/2)).$$

La costruzione delle bitangenti che lasciano i centri in semipiani diversi è un esempio classico che ricordiamo qui per convenienza del lettore.

La proprietà alla base della costruzione segue dalla similitudine dei triangoli $OPQ$ e $OP'Q$ in Figura 8, da cui
$$\frac{OQ}{O'Q}=\frac{OP}{O'P'}=\frac{r}{R}.$$
Non resta dunque che costruire prima il punto $Q$ con tale proprietà e poi le tangenti alle 2 circonferenze passanti per $Q$.
\begin{center}
	\hfill
	\psset{unit=10mm}
	\pspicture(-2,-2.4)(3,6)
	\rput(0.5,-2){Figura 5}
	\psframe[linewidth=.6\pslinewidth](-2,-1.5)(3,6)
	
	\pstGeonode[PointSymbol={*,*,none},PosAngle=180](0,0){O}(0,4){O'}(1,0){r}(1.5,4){R}(2.5,0){R'}
	\pstCircleOA{O}{r}
	\pstCircleOA{O'}{R}
	\pstRotation[RotAngle=-10,PosAngle=-70]{O}{r}[A]
	\pstRotation[RotAngle=-10,PosAngle=-90]{O}{R'}[B]
	\pstLineAB[linecolor=blue]{O}{A}
	\Aput[2pt]{$r$}
	\pstLineAB[linecolor=blue]{A}{B}
	\Aput[2pt]{$r'$}
	\pstLineAB[linecolor=blue]{O'}{R}
	\Aput[2pt]{$r'$}
	
	\endpspicture
	\hfill
	\psset{unit=10mm}
	\pspicture(-2,-2.4)(3,6)
	\rput(0.5,-2){Figura 6}
	\psframe[linewidth=.6\pslinewidth](-2,-1.5)(3,6)
	
	\pstGeonode[PointSymbol={*,*,none},PosAngle=180](0,0){O}(0,4){O'}(1,0){r}(1.5,4){R}(2.5,0){R'}
	\pstCircleOA{O}{r}
	\pstCircleOA{O'}{R}
	\pstRotation[RotAngle=-10,PosAngle=-90]{O}{r}[A]
	\pstRotation[RotAngle=-10,PosAngle=-110]{O}{R'}[B]
	\pstLineAB[linecolor=blue]{O}{A}
	\Aput[2pt]{$r$}
	\pstLineAB[linecolor=blue]{A}{B}
	\Aput[2pt]{$r'$}
	
	\pstTranslation[PointSymbol=none]{B}{O'}{A}[A']
	\pstInterLL{A}{A'}{O}{O'}{Q}
	
	\pstLineAB[nodesep=-1,linecolor=blue]{B}{O'}
	\pstLineAB[nodesep=-1,linecolor=blue]{A}{Q}
	\pstLineAB[nodesep=-0.8,linecolor=blue]{O}{O'}
	
	\endpspicture
	\hfill\mbox{}
	
	\hfill\psset{unit=10mm}
	\pspicture(-2,-2.4)(2,6)
	\rput(0,-2){Figura 7}
	\psframe[linewidth=.6\pslinewidth](-2,-1.5)(2,6)
	
	\pstGeonode[PointSymbol={*,*,none},PosAngle={-45,45}](0,0){O}(0,4){O'}(1,0){r}(1.5,4){R}(2.5,0){R'}
	\pstCircleOA{O}{r}
	\pstCircleOA{O'}{R}
	\pstRotation[RotAngle=-10,PointSymbol=none]{O}{r}[A]
	\pstRotation[RotAngle=-10,PointSymbol=none]{O}{R'}[B]
	\pstTranslation[PointSymbol=none]{B}{O'}{A}[A']
	\pstInterLL[PosAngle=50]{A}{A'}{O}{O'}{Q}
	\pstCircleAB[linecolor=green]{O}{Q}
	\pstCircleAB[linecolor=green]{O'}{Q}
	\pstMiddleAB[PointSymbol=none]{O}{Q}{O''}
	\pstMiddleAB[PointSymbol=none]{O'}{Q}{O'''}
	\pstInterCC[PointName=none]{O}{r}{O''}{Q}{P}{S}
	\pstInterCC[PointName=none]{O'}{R}{O'''}{Q}{P'}{S'}
	\pstLineAB[nodesep=-0.8,linecolor=blue]{O}{O'}
	
	\endpspicture
	\hfill
	\psset{unit=10mm}
	\pspicture(-2,-2.4)(2,6)
	\rput(0,-2){Figura 8}
	\psframe[linewidth=.6\pslinewidth](-2,-1.5)(2,6)
	
	\pstGeonode[PointSymbol={*,*,none},PosAngle={-45,45}](0,0){O}(0,4){O'}(1,0){r}(1.5,4){R}(2.5,0){R'}
	\pstCircleOA{O}{r}
	\pstCircleOA{O'}{R}
	\pstRotation[RotAngle=-10,PointSymbol=none]{O}{r}[A]
	\pstRotation[RotAngle=-10,PointSymbol=none]{O}{R'}[B]
	\pstTranslation[PointSymbol=none]{B}{O'}{A}[A']
	\pstInterLL{A}{A'}{O}{O'}{Q}
	\pstCircleAB[linecolor=green]{O}{Q}
	\pstCircleAB[linecolor=green]{O'}{Q}
	\pstMiddleAB[PointSymbol=none]{O}{Q}{O''}
	\pstMiddleAB[PointSymbol=none]{O'}{Q}{O'''}
	\pstInterCC[PointNameB=none,PosAngle=180]{O}{r}{O''}{Q}{P}{S}
	\pstInterCC[PointNameB=none]{O'}{R}{O'''}{Q}{P'}{S'}
	
	\pstLineAB[nodesep=-0.8,linecolor=blue]{O}{O'}
	\pstLineAB[linecolor=blue]{O'}{P'}
	\pstLineAB[linecolor=blue]{O}{P}
	\pstRightAngle[linecolor=blue]{O'}{P'}{Q}
	\pstRightAngle[linecolor=blue]{O}{P}{Q}
	
	\pstLineAB[nodesep=-1,linewidth=1.5\pslinewidth]{P}{P'}
	\pstLineAB[nodesep=-1,linewidth=1.5\pslinewidth]{S}{S'}
	
	\endpspicture
	\hfill\mbox{}
\end{center}

La costruzione è allora la seguente.
\begin{itemize}
	\item  Tracciamo una semiretta uscente da $O$, diversa da $OO'$, sulla quale segniamo 2 punti $A$ e $B$ tali che $OA=r$ e $AB=r'$ (Figura 5).

	\item  Tracciamo la parallela ad $O'B$ passante per $A$. Tale retta incontra $OO'$ in un punto $Q$. Per il teorema di Talete si ha che
	$$\frac{OQ}{QO'}=\frac{OA}{AB}=\frac{r}{r'}$$
	e pertanto $Q$ è il punto in cui le bitangenti che ci interessano incontrano il segmento $OO'$ (Figura 6).

	\item Tracciamo le 2 circonferenze di diametro $OQ$ e $O'Q$, rispettivamente (Figura 7).  Le intersezioni tra queste 2 circonferenze e le 2 circonferenze iniziali individuano i punti di tangenza delle bitangenti cercate (Figura 8).
\end{itemize}

}

\renewcommand{\SoluzioniAlternative}{
}

\renewcommand{\Osservazioni}{
}

\renewcommand{\Commenti}{
}